\section{Surveillance over time and against strategic traders}\label{app:dynamic}

The cost results treat the population as fixed. Real populations drift: algorithms enter, exit and are retuned, some of them in response to the supervision itself, and the groups being watched choose how to trade. This section gives three results: a speed limit on tracking a moving population, a floor on what probing can learn about traders who learn from prices, and an equilibrium in which surveillance without tags pushes algorithms together and never detects them.

\subsection{Tracking a moving population}

Fix the setting of Theorem~\ref{thm:kcohort}: $K\ge2$ groups with rates $r_c+\Delta\eta\,s_i$, rescaled nodes $s$ and weights $\omega$ in a rich class, persistence and impact known, and the market Schur stable on the class. Time runs in periods. In period $t$ the population is $\theta_t=(s_t,\omega_t)$, and the supervisor runs $n$ reset experiments. A path is \emph{admissible at speed $v$} if every $\theta_t$ lies in the class and $\|\theta_{t+1}-\theta_t\|_\infty\le v$. At date $T$ the supervisor estimates $\omega_T$ from all data up to $T$.

\begin{theorem}[Tracking a moving population]\label{thm:track}
There are constants $c,C,\bar v>0$ and $v^*_0>0$, depending only on the class, such that, with $\Gamma=\sigma^2v(\Delta\eta)^{-(4K-2)}/n$:
\begin{enumerate}[label=(\roman*),itemsep=1pt]
\item \emph{Lower bound.} For every estimator and every $T$ large enough, some path admissible at speed $v\le\bar v$ has $\Pr\bigl(\max_i|\hat\omega_{T,i}-\omega_{T,i}|\ge c\min(1,\Gamma^{1/3})\bigr)\ge1/4$.
\item \emph{Upper bound.} If $\Gamma\le1/C$, Prony's method applied to the taps averaged over the last $m\asymp(\sigma^2(\Delta\eta)^{-(4K-2)}/(nv^2))^{1/3}$ periods has $\max_i|\hat\omega_{T,i}-\omega_{T,i}|=O_P(\Gamma^{1/3})$ along every path admissible at speed $v$, for $T\ge m$.
\item \emph{Speed limit.} If $v\ge v^*:=v^*_0\,n(\Delta\eta)^{4K-2}/\sigma^2$, no estimator tracks the weights to accuracy below $c$, however long the record.
\item \emph{Market-wide questions.} For a linear question $F_g$ with square-summable coefficients whose answer moves by at most $v_F$ per period along the path, the window average of the truncated estimator of Theorem~\ref{thm:questions}(i) tracks it with error $O_P((\sigma^2v_F/n)^{1/3})$, whatever $K$ and $\Delta\eta$.
\end{enumerate}
\end{theorem}

\begin{proof}
(i) Use the family $R_c$ of Lemma~\ref{lem:everyK}: for $c$ near zero, the $K$-node rules whose nodes are the zeros of $\ell_K+c\,\ell_{K-1}$ match all moments of orders up to $2K-2$, so all members give populations that agree in the first $2K-1$ rescaled moments. Their nodes are simple zeros of a polynomial that is analytic in $c$, so nodes and weights are analytic in $c$; by Lemma~\ref{lem:everyK}(iii) the weights are not all constant, so some weight $\omega_{i_0}(c)$ has $|\omega_{i_0}'(c^*)|\ge2b>0$ at some $c^*$ arbitrarily close to zero, and $|\omega_{i_0}'|\ge b$ on an interval $I=[c^*,c^*+\iota]$ on which the rules lie in the class (scaled to total weight $W$). Let $B=\sup_I\|\tfrac{d}{dc}(s,\omega)\|_\infty$. On $I$, the rescaled moments of orders $p\ge2K-1$ are Lipschitz in $c$ with constant at most $C_1(1+p)$, so, exactly as in the proof of Theorem~\ref{thm:kcohort}, the price responses of the populations at $c$ and $c'$ differ in $\ell^2$ by at most $C_2|c-c'|(\Delta\eta)^{2K-1}$, and one reset experiment carries divergence at most $C_3(c-c')^2(\Delta\eta)^{4K-2}/\sigma^2$.

Compare two paths. On the first, $c_t=c^*$ for all $t$. On the second, $c_t=c^*$ up to date $T-m$ and $c_t=c^*+(v/B)(t-T+m)$ afterwards; it is admissible at speed $v$ as long as $vm/B\le\iota$. The two data laws differ only in the last $m$ periods, in the $j$-th of which the populations are $vj/B$ apart along the family, so the divergence between them is at most $\sum_{j=1}^mnC_3(vj/B)^2(\Delta\eta)^{4K-2}/\sigma^2\le nC_3v^2m^3(\Delta\eta)^{4K-2}/(B^2\sigma^2)$, using $\sum_{j\le m}j^2\le m^3$. At date $T$ the weights $\omega_{i_0}$ differ by at least $bvm/B$. Choose $m=\min\{\lfloor(B^2\sigma^2/(2nC_3v^2(\Delta\eta)^{4K-2}))^{1/3}\rfloor,\lfloor\iota B/v\rfloor\}$, which is at least one for $v\le\bar v$ and small $\Delta\eta$. The divergence is then at most $1/2$, and Le Cam's two-point argument gives an error of at least half the separation, $bvm/(2B)$, with probability at least $1/4$ under one of the two paths. If the first term attains the minimum, the separation is of order $\Gamma^{1/3}$; if the second does, it is $b\iota/2$, a constant.

(ii) The rescaled moments of a population are linear in its merged measure, and on the class the rescaled moment of order $p$ is Lipschitz in $(s,\omega)$ with a constant that does not depend on $\Delta\eta$, because $|s_i|\le1$. So the average of the rescaled moments over the last $m$ periods differs from those at date $T$ by at most $C_4vm$. The averaged taps estimate the rescaled moments of that average with error $O_P(\sigma(\Delta\eta)^{-(2K-1)}(nm)^{-1/2})$, as in Theorem~\ref{prop:esprit}. The Prony map has a Lipschitz inverse near the class (Theorem~\ref{prop:esprit}), so the weight error is $O_P(vm+\sigma(\Delta\eta)^{-(2K-1)}(nm)^{-1/2})$, and the stated $m$ balances the two terms at order $\Gamma^{1/3}$. The condition $\Gamma\le1/C$ keeps the moment error within the radius of the inverse.

(iii) This is the second case of (i): if $v\ge v^*$ with $v^*_0$ small enough, the window $\iota B/v$ is the binding one, the divergence is at most $1/2$, and the separation is a constant.

(iv) The bias of the window average is at most $v_Fm$ plus the truncation error, and its noise is $O_P(\sigma\sqrt{A_g}(nm)^{-1/2})$; balancing gives the rate, with no dependence on the population beyond the bound on the weights.
\end{proof}

\paragraph{What this means.} Theorem~\ref{thm:kcohort} says composition is expensive to learn; Theorem~\ref{thm:track} says that when composition moves, the expense becomes a speed limit. The supervisor's error falls only like the cube root of the data rate, and populations that change faster than $v^*\propto n(\Delta\eta)^{4K-2}/\sigma^2$ cannot be tracked at all. In the proof, a population that drifts across the distance between two members of the hard family in a few identification times cannot be told from one that stays still. At illustrative parameters (memories with half-lives of 5 to 9 days), where telling such a pair apart takes thousands of years of daily prices, composition is untrackable unless it is stable over millennia. Market-wide questions are tracked at the cube-root rate whatever the number of groups; and along the hard directions their answers hardly move, because those directions preserve the low moments that market-wide answers mostly depend on.

\paragraph{Numerical check.} With two groups moving along the family and the window chosen to minimize the median error, the weight error grows with the speed at fitted slope $0.33$ (against $1/3$) and with the spacing at slope $-2.01$ (against $-2$, the cube root of $(\Delta\eta)^{-6}$). The error for the discounted feedback grows with speed at slope $0.34$; its dependence on the spacing, slope $0.98$, comes only through the speed of its answer, which along the family is of order $(\Delta\eta)^3$ (\nolinkurl{research_extension/tracking/tracking_check.py}).

\subsection{Traders who learn from prices: an observer effect}

Suppose the traders adjust their rules in response to what they earn, and so in response to the supervisor's probes. A probe of size $A$ moves a group's position by an amount proportional to $A$ and the price by an amount proportional to $A$, so its profit and loss is proportional to $A^2$: with position $w_im_t$ and price changes $\Delta x_{t+1}$, the profit over a probe's response is $w_iA^2\Pi_i$, where $\Pi_i=\sum_tm_t^{(1)}\Delta x_{t+1}^{(1)}$ is computed on the unit response. A rule that moves its weight by $\eta$ times realized profit therefore moves by at most $\eta A^2\max_i|w_i\Pi_i|$ per probe. Call a class of learning rules \emph{of sensitivity $\gamma$} if, with $n$ probes of size $A$ per period, the population can move at speed up to $v_0+\gamma nA^2$, and the supervisor knows only this bound.

\begin{proposition}[An observer effect]\label{prop:observer}
In the setting of Theorem~\ref{thm:track}, with $n$ probes of size $A$ per period and learning rules of sensitivity $\gamma$, every probing policy leaves a minimax tracking error of at least
\[
 c\min\Bigl(1,\bigl(\sigma^2(\Delta\eta)^{-(4K-2)}(v_0/(nA^2)+\gamma)\bigr)^{1/3}\Bigr)\ \ge\ c\min\Bigl(1,\bigl(\sigma^2\gamma(\Delta\eta)^{-(4K-2)}\bigr)^{1/3}\Bigr).
\]
For a market-wide question whose answer moves at most $\gamma_F$ per unit of probe energy, the floor is $c\min(1,(\sigma^2\gamma_F)^{1/3})$.
\end{proposition}

\begin{proof}
A probe of size $A$ is a unit probe with noise $\sigma/A$, so Theorem~\ref{thm:track}(i) applies with $\sigma$ replaced by $\sigma/A$ and $v$ by $v_0+\gamma nA^2$: then $\Gamma=\sigma^2(v_0+\gamma nA^2)(\Delta\eta)^{-(4K-2)}/(nA^2)$, which is at least $\sigma^2\gamma(\Delta\eta)^{-(4K-2)}$. The lower bound of Theorem~\ref{thm:track}(i) needs only that some admissible path moves along the hard family at that speed, which a class of learning rules known only through its sensitivity allows. The market-wide case is the same computation with Theorem~\ref{thm:track}(iv).
\end{proof}

Probing harder buys information at rate $A^2$ and moves the target at rate $A^2$, so the two cancel: surveillance of traders who learn from prices has a floor that no probing design removes. Two caveats bound the result. It is a minimax statement over learning rules of a given sensitivity; a supervisor who knows the learning rule can anticipate the drift it causes. And the floor for composition carries the factor $(\Delta\eta)^{-(4K-2)}$, while the floor for market-wide questions does not.

\subsection{A surveillance and hiding game}

A group that expects to be targeted once identified may choose to trade like its neighbours. Consider two untagged groups and let one of them choose the spacing $D\in(0,D_0]$ between its memory and its neighbour's, where $D_0$ is the spacing of its preferred strategy and moving away from it costs $k(D_0-D)^2$. A supervisor with $N$ observations runs a test with false-alarm probability $\alpha$; by Theorems~\ref{thm:kcohort} and~\ref{prop:passive}, its detection probability depends on the data and the spacing through $ND^6$ at leading order, $\pi(ND^6)$ with $\pi$ continuous and nondecreasing and $\pi(0)=\alpha$. Detection costs the group $L>0$. Tags, if the supervisor imposes them, make the group's weight identifiable at a cost that does not depend on $D$ (Theorem~\ref{thm:partial}).

\begin{proposition}[Surveillance and hiding]\label{prop:hidegame}
Let $D^*(N)$ minimize $k(D_0-D)^2+L\pi(ND^6)$ over $(0,D_0]$.
\begin{enumerate}[label=(\roman*),itemsep=1pt]
\item \emph{Crowding.} If $\pi$ is differentiable at $ND_0^6$ with positive derivative, then $D^*(N)<D_0$.
\item \emph{Surveillance without tags is self-limiting.} $\limsup_{N\to\infty}\pi(ND^*(N)^6)\le\alpha+kD_0^2/L$. When the stakes are high, $L\gg kD_0^2$, the group is detected about as often as a false alarm, however much data the supervisor has.
\item \emph{Tags remove crowding.} If the detection probability does not depend on $D$, then $D^*=D_0$.
\end{enumerate}
\end{proposition}

\begin{proof}
(i) The derivative of the objective at $D_0$ is $6LND_0^5\pi'(ND_0^6)>0$, so a small move below $D_0$ lowers it. (ii) For any $c>0$ and $N$ large enough that $cN^{-1/6}\le D_0$, optimality against $D=cN^{-1/6}$ gives $L\pi(ND^{*6})\le k(D_0-cN^{-1/6})^2+L\pi(c^6)\le kD_0^2+L\pi(c^6)$. Letting $c\to0$ and using continuity of $\pi$ at zero gives the bound. (iii) The cost $k(D_0-D)^2$ is minimized at $D_0$.
\end{proof}

The proposition gives the paper's results on cost an equilibrium reading. More data without tags make high-stakes algorithms trade alike, and so make composition harder to learn, while the detection probability stays near the false-alarm level. Tags undo the crowding. Two predictions follow that algorithm-level data could test: where algorithm identification is introduced, the dispersion of memory rates across algorithm classes should rise; and in untagged markets, the memories of strategies with the most to lose from targeting should converge as surveillance improves. With the Gaussian test, the optimal spacing in the high-stakes case falls like $N^{-1/4}$ and detection stays at the $5\%$ false-alarm level up to $N=10^{16}$; in the low-stakes case the group hides for moderate $N$ and stops hiding once data are plentiful (\nolinkurl{research_extension/tracking/hiding_game_check.py}). The model has one strategic group and one neighbour; with several groups each hides from its nearest neighbour, and the equilibrium of that game is not studied here.

\paragraph{Machine check.} The divergence of a moving window, Le Cam's two-point step, the bias--noise trade-off, the observer-effect inequality and parts (ii) and (iii) of Proposition~\ref{prop:hidegame} are machine-checked (\texttt{p1\_track\_kl\_window}, \texttt{p1\_track\_two\_point}, \texttt{p1\_track\_tradeoff}, \texttt{p1\_observer\_floor}, \texttt{p1\_hiding\_detection\_bound}, \texttt{p1\_hiding\_tag\_no\_crowding}, \nolinkurl{Paper1TrackingGame.lean}). The analytic properties of the family $R_c$, the moment Lipschitz bounds and the closed-loop step are on paper.
